\section{Introduction}\label{sec:intro}

The final layer of a classifier is stochastic. The softmax defines a categorical distribution over classes, and at inference we either sample from it or take its \texttt{argmax}; at training time the same distribution is what the cross-entropy loss scores. The hidden layers, by contrast, are deterministic: a ReLU activation passes its exact value forward, unchanged from pass to pass. This asymmetry is so familiar that it is rarely questioned. We ask the sharp version of the question it invites: {\bf what if we apply the output-layer operation at an intermediate layer?} That is, form a probability distribution over a hidden layer's activations and draw a \emph{sample} from it, then pass that sample forward in place of the deterministic activations. The literal form of this idea treats a hidden layer exactly like the output layer --- a softmax over its units gives a categorical distribution, and we sample one unit, just as we sample one class.

Hidden-layer stochasticity is not new; it is the shared mechanism behind a remarkable range of methods. Variational autoencoders sample a continuous latent through the reparameterization trick \citep{kingma2014vae,rezende2014stochastic}; Gumbel-Softmax and the Concrete distribution sample discrete latents differentiably \citep{jang2017gumbel,maddison2017concrete}; stochastic neurons pass binary samples with straight-through gradients \citep{bengio2013stochastic}; dropout injects multiplicative Bernoulli noise that doubles as a Bayesian approximation \citep{srivastava2014dropout,gal2016dropout}; and the deep variational information bottleneck deliberately makes a representation stochastic to limit the information it carries \citep{alemi2017deepvib,tishby2015deep}. In every case the stochasticity is introduced \emph{for a purpose} --- generative modeling, discrete latents, regularization, compression --- and \emph{at a fixed place} chosen for that purpose. A practitioner who simply wants to know the cost (in accuracy) and the benefit (in calibration, generalization, and robustness) of sampling a hidden layer, and where in the network to place it, has no controlled answer to point to.

\para{The gap.} No prior work isolates intermediate-layer activation sampling as the variable of study itself. The methods above each entangle the sampling operation with an objective, an architecture, and a placement, so the effect of the sampling \emph{per se} is never measured against an otherwise-identical deterministic network. The softmax-sampling analogy --- treat a hidden layer exactly as the output layer, softmax over its units to a categorical distribution and sample one --- is almost never made explicit, and its consequences are essentially unstudied. We close this gap with a single drop-in module and a controlled comparison.

\para{Our approach.} We build one \samplinglayer{} that, placed after a hidden ReLU, turns the incoming activations into a distribution and draws a sample. Holding architecture, data, and optimizer fixed, we compare seven activation operators: a \detmode{} reference; two noise baselines (\dropout{} and additive \gaussnoise); and the three faces of the idea the literature identifies --- continuous \gaussreparam{} (VAE-style), binary Bernoulli straight-through (\bernoullist{}, the Bengio stochastic neuron), and categorical Gumbel-Softmax in its soft (\gumbelsoft) and hard one-hot (\gumbelst) forms, the direct softmax analogue that collapses the layer to a single sampled unit. We evaluate on \mnist{} (MLP) and \cifar{} (CNN) across 3 seeds --- 117 trained models in total --- measuring accuracy, \nll{}, \ece{} calibration, the train--test generalization gap, gradient variance at the sampling layer, and single-sample versus MC-averaged inference, with sweeps over temperature/noise scale, layer position, and input-noise robustness.

\para{What we find.} Three results stand out. First, {\bf the distribution type dominates everything.} Continuous \gaussreparam{} sampling is essentially free --- $97.7\%$ vs $97.8\%$ deterministic accuracy on \mnist{} (paired $t$-test $p\!=\!0.83$, not significant) --- and on the harder \cifar{} it \emph{beats} the deterministic reference, cutting \ece{} from $0.112$ to $0.074$ and \nll{} from $0.94$ to $0.78$ at equal-or-better accuracy. The literal categorical analogue does the opposite: it is the most expensive operator on \mnist{} ($-10.4$ points for \gumbelst, $p\!=\!0.048$) and \emph{collapses to chance} ($10\%$) on \cifar{}. Second, {\bf sampling regularizes.} Every stochastic variant shrinks the train--test gap and most lower \ece{} --- the benefit the literature predicts --- so the gentle variants trade a tiny (or zero) accuracy cost for better-behaved probabilities. Third, {\bf MC-averaging recovers the single-sample cost}, exactly as averaging softmax samples approximates the class expectation: dropout climbs from $96.6\%$ to $97.9\%$ as the number of test samples $K$ grows from $1$ to $50$. Underlying all three, the gradient-variance ranking of the operators (lowest for Bernoulli straight-through, highest for hard categorical) tracks the accuracy cost almost exactly, giving a mechanistic explanation for why categorical sampling is so hard.

\para{Contributions.}
\begin{itemize}[leftmargin=*,itemsep=0pt,topsep=0pt]
  \item We propose \samplinglayer, a single drop-in module that turns any hidden layer into a distribution and samples from it, exposing seven activation operators --- a deterministic reference, two noise baselines, and the continuous, binary, and categorical faces of activation sampling --- behind one interface.
  \item We conduct a controlled comparison that isolates intermediate-layer activation sampling as the variable of study, holding architecture, data, and optimizer fixed across \mnist{} and \cifar{}, 3 seeds, and 117 models, with accuracy, \nll, \ece, generalization gap, gradient variance, MC-averaging, and sweeps over temperature, layer position, and input noise.
  \item We show that continuous \gaussreparam{} sampling is free and even beneficial ($-0.03$ pts on \mnist; $+0.5$ pts accuracy and $0.74$ vs $0.94$ \nll{} on \cifar) while the direct categorical softmax analogue is the most expensive variant and collapses to chance on \cifar, making explicit that ``sampling a hidden layer like the softmax'' is distribution-type-dependent rather than a free drop-in.
  \item We explain the cost mechanistically: the gradient variance induced at the sampling layer ranks the operators in the same order as their accuracy cost, identifying high-variance gradient estimation as the bottleneck for categorical sampling.
\end{itemize}

The rest of the paper is organized as follows. \Secref{sec:related} places the seven operators within prior work on stochastic representations; \Secref{sec:method} details the \samplinglayer{} and experimental protocol; \Secref{sec:results} reports the core comparison, sweeps, and the \cifar{} confirmation; and \Secref{sec:discussion} and \Secref{sec:conclusion} interpret the findings and outline next steps.
